A spacecraft of mass $m = 5.0 \times 10^4$\,kg approaches in a parabolic orbit
$AB$, with respect to Mars. When the spacecraft reaches point $B$ of least
distance to the center of Mars, $r_B = 6.8 \times 10^6$\,m, it undergoes an
instantaneous deceleration using its rockets and goes into a perfectly
calculated orbit so that it will touch the Martian surface exactly at point
$C$, diametrically opposite $B$, as shown in the figure.

\ioaafig{2021_TH_Q3-f1.pdf}

\begin{description}
\item[(3.1)] Determine the speed ($\mathrm{km\,s}^{-1}$) of the spacecraft at
  point $B$ just before the deceleration. \hfill[3.0\,pt]
\item[(3.2)] Calculate the total energy ($J$) of the spacecraft as it is
  moving between points B and C. \hfill[4.0\,pt]
\item[(3.3)] Calculate the speed ($\mathrm{km\,s}^{-1}$) of the spacecraft at
  point $C$. \hfill[3.0\,pt]
\end{description}
